\documentclass[10pt,a4paper]{amsart}
\usepackage[margin=19mm]{geometry}
\usepackage{amsmath,amssymb,mathtools}
\usepackage[scr=rsfs,cal=euler]{mathalfa}
\usepackage{xfrac}
\usepackage[unicode,hidelinks,bookmarks=false]{hyperref}
\newcommand{\ie}{i.e.\ }
\newcommand{\cf}{cf.\ }
\newcommand{\resp}{respectively\ }
\newcommand{\Subsection}{\subsection}
\providecommand{\nobreakdash}{\nobreak\hbox{-}}
\setlength{\emergencystretch}{2em}
\multlinegap0pt
\usepackage{bm}
\usepackage{bbm}
\usepackage{dsfont}
\usepackage{xspace}
\usepackage{cancel}
\usepackage{mathtools}
\usepackage{amsthm}
\makeatletter
\@ifundefined{theorem}{\newtheorem{theorem}{Theorem}}{}
\makeatother
\newcommand{\Ra}{\Rightarrow}
\newcommand{\SM}{{\setminus}}
\newcommand{\w}{\omega}
\title[On $\kappa$-Frechet-Urysohn topological groups]{On $\kappa$-Frechet-Urysohn topological groups}
\author{Saak Gabriyelyan, Alexander V. Osipov, Evgenii Reznichenko}
\date{Source: arXiv:2602.04647v2 (CC BY 4.0)}
\begin{document}
\maketitle

\noindent\emph{Source and licence.} On $\kappa$-Frechet-Urysohn topological groups. Version arXiv:2602.04647v2 (CC BY 4.0), distributed under the Creative Commons Attribution 4.0 International licence, as recorded in the arXiv licence metadata for this exact version and shown on the abstract page https://arxiv.org/abs/2602.04647v2. Full rights notice: licenses/arxiv-2602.04647-CC-BY-4.0.txt.\par\smallskip

\noindent\textit{Editorial scope.} Counted unit: the Theorem whose source label is \texttt{t:kFU-group-charac} in the article (source0.tex lines 363--370, source label \texttt{t:kFU-group-charac}), delivered together with the article's own complete original proof of it (source0.tex lines 372--392, one \texttt{proof} environment of 21 lines). No mathematics is written, repaired or re-derived: statement and proof are the article's own text line for line. Required context retained uncounted: none. Every definition, display and label of those lines is retained unchanged. Omitted from this package and not redistributed: 1--362, 371--371, 393--598. Nothing inside a retained excerpt is omitted or altered: every retained body is delivered exactly as the article's lines are written, so no omission receipt of any kind accompanies this package. Citations: the retained excerpts carry no citation command at all, so no bibliography entry of the article is transcribed and no reference is redistributed. Reference adaptation: none was needed and none was made; every reference inside the delivered unit resolves inside it, so no unresolved reference or citation is produced. Printed slips kept literal: no printed word, symbol, hypothesis or number of the counted statement or of its proof was altered. Layout and class: the article's own class is \texttt{elsarticle} with packages including amsfonts, amsmath, amscd, amssymb, pb-diagram, xy, amsthm, color; the wrapper is the lane's pinned \texttt{amsart} editorial wrapper at 10pt on A4, which restates the theorem-like environments the retained text needs and adds the article's own macro definitions that the retained text needs (\texttt{\textbackslash Ra}, \texttt{\textbackslash SM}, \texttt{\textbackslash w}), with extra wrapper packages (bm, bbm, dsfont, xspace, cancel, mathtools). The delivered numbering is the unit's own. Assets: no retained span contains a figure environment, a graphics-inclusion command, a PDF-page inclusion, a table or a TikZ picture; no image file is copied into this package, the package declares no asset dependency and no image file is redistributed with it. Licence: version arXiv:2602.04647v2 of this article is distributed under the Creative Commons Attribution 4.0 International licence, as recorded in the arXiv licence metadata for this exact version and shown on the abstract page https://arxiv.org/abs/2602.04647v2. A full rights notice is delivered at licenses/arxiv-2602.04647-CC-BY-4.0.txt. Characters: every retained excerpt is pure ASCII, so the delivered engine, XeLaTeX with its default Latin Modern text font, reports no missing character.
\begin{theorem} \label{t:kFU-group-charac}
For a topological group $G$, the following assertions are equivalent:
\begin{enumerate}
\item[{\rm(i)}] for every sequence $\{\mathcal{V}_n: n\in \w\}$ of countable families $\mathcal{V}_n=\{V^i_n: i\in \w\}$ of open subsets of $G$ such that $\mathcal{V}_n$ weakly converges to $e$ for every $n\in\w$, there are strictly increasing sequences  $(i_k)$ and $(n_k)$ in $\w$ and a sequence $\{x_k\}_{k\in\w}$ such that $x_k\to e$ and $x_k\in V^{i_k}_{n_k}$ for every $k\in\w$;
\item[{\rm(ii)}] for every sequence $\{U_n\}_{n\in\w}$ of open subsets of $G$ such that $e\in \overline{U_n}$ for every $n\in\w$, there are a strictly increasing sequence $(n_k)\subseteq \w$ and a sequence $\{x_k\}_{k\in\w}$ in $G$ such that $x_k\in U_{n_k}$ $(k\in\w)$ and $x_k\to e$;
\item[{\rm(iii)}] $G$ is a $\kappa$-Fr\'{e}chet--Urysohn space.
\end{enumerate}
\end{theorem}

\begin{proof}
(i)$\Ra$(ii) For every $i\in \w$, we set $V^i_n:=U_n$. Now (ii) immediately follows from (i).
\smallskip

(ii)$\Ra$(iii) Let $U$ be an open subset of $G$ such that $e\in \overline{U}\SM U$. For every $n\in\w$, set $U_n:=U$. Then (iii) immediately follows from (ii).
\smallskip

(iii)$\Ra$(i) Let $\{\mathcal{V}_n: n\in \w\}$  be a family satisfying (i). Choose an arbitrary one-to-one sequence $\{p_n\}_{n\in\w}$ in $G\SM \{e\}$  converging to $e$ (such a sequence exists by the $\kappa$-Fr\'{e}chet--Urysohness applying to the open set $G\SM\{e\}$). For every $n\in\w$, set $p_n\cdot\mathcal{V}_n:=\{p_n V^i_n: i\in \w\}$. Choose a sequence $\{W_n: n\in \w\}$ of open neighborhoods of $e$ such that
\begin{equation} \label{equ:kFU-group-charac-1}
%W_{n+1}+W_{n+1}\subseteq W_n, \;\; (p_n+W_n) \cap (p_{n+1}+W_{n+1})=\emptyset \;\; \mbox{ and }\;\;
e\not\in \overline{p_n W_n}
\end{equation}
for every $n\in \w$. For each $0\leq n\leq i$, set
\begin{equation} \label{equ:kFU-group-charac-2}
\mathcal{O}^i_n=p_nV^i_n\cap p_nW_n.
\end{equation}
We claim that $e\in \overline{\bigcup \{\mathcal{O}^i_n: 0\leq n\leq i\}}$.  Indeed, let $U$ be a neighborhood of $e$. Choose a neighborhood $U_1$ of $e$ such that $U_1\cdot U_1\subseteq U$. Since $p_n\to e$, there is $m\in\w$ such that $p_n\in U_1$ for each $n\geq m$.  As $\mathcal{V}_m$ weakly converges to $e$, there are $j\geq m $ and a point $s_m^j\in V_m^j$ such that $s_m^j\in U_1\cap W_m$. Then $p_m\cdot s_m^j\in U_1\cdot U_1\subseteq U$ and $p_m\cdot s_m^j\in \mathcal{O}^j_m$.
\smallskip

Since $G$ is $\kappa$-Fr\'{e}chet--Urysohn, there is a sequence $(z_k)\subseteq \bigcup \{\mathcal{O}^i_n: 0\leq n\leq i\}$ converging to $e$. For every $k\in\w$, choose indices  $i_k$ and $n_k$ such that $z_k\in \mathcal{O}^{i_k}_{n_k}$. The sequence $(n_k)$ cannot be bounded because of (\ref{equ:kFU-group-charac-1}) and the inclusion $\mathcal{O}^{i_k}_{n_k} \subseteq p_{n_k}W_{n_k}$. Taking into account that $j_k\geq n_k$ it follows that also $(j_k)$ is unbounded. Passing to subsequences if needed, we can assume that $(j_k)$ and $(n_k)$ are strictly increasing.  For every $k\in\w$, set $x_k:= p_{n_k}^{-1}\cdot z_k $. Then, by (\ref{equ:kFU-group-charac-2}), $x_k\in V^{i_k}_{n_k}$ and, clearly, $x_k\to e$. Thus (i) is satisfied.
\end{proof}

\end{document}
